% BLOCK13_STAGE2_MANAGED
% Lecture 13: type deduction, lambdas, std::function and std::optional
\section{Лекция 13. Вывод типов, лямбды и вызываемые объекты}
\label{sec:lecture-13}

В предыдущих блоках мы всё чаще встречали длинные типы:
\texttt{std::unique\_ptr<Product>}, iterator types, template arguments и callable
objects. Современный C++ позволяет не дублировать тип там, где compiler уже
может его вывести. Но \texttt{auto} не означает «динамический тип», а lambda
не «анонимная функция» в смысле обычного function pointer.

В этой лекции мы строим прикладную линию:
\[
\text{type deduction}
\rightarrow
\text{lambda}
\rightarrow
\text{callable API}
\rightarrow
\texttt{std::function}
\rightarrow
\texttt{std::optional}.
\]

Итоговый design-вопрос звучит так: как написать API, который принимает
поведение и возвращает отсутствие результата без magic values, оставаясь
понятным по типам?

\begin{importantbox}[Сквозная идея]
Вывод типа убирает повторение, но не отменяет type system. Lambda создаёт object
уникального closure type. \texttt{std::function} стирает конкретный callable
type ради единого runtime-интерфейса. \texttt{std::optional<T>} выражает
«либо T, либо отсутствует значение».
\end{importantbox}

\subsection{\texttt{auto}: compiler выводит тип из initializer}

Простейший пример:

\begin{cppcode}
auto count = 42;
auto ratio = 0.5;
auto name = std::string{"C++"};
\end{cppcode}

Здесь types соответственно выводятся из initializer. После deduction variable
имеет обычный статический C++ type. Никакой runtime-перемены типа не появляется.

\begin{whybox}[Зачем использовать auto, если тип и так известен?]
Главная польза не в экономии символов, а в устранении дублирования. Если
expression уже точно задаёт тип, повторять его слева иногда только создаёт место
для рассинхронизации.
\end{whybox}

Например:

\begin{cppcode}
std::vector<std::string> names{"Ada", "Bjarne"};
auto it = names.begin();
\end{cppcode}

Тип iterator может быть длинным и зависит от container. \texttt{auto} сохраняет
важную информацию: \texttt{it} имеет именно тот type, который возвращает
\texttt{begin()}.

\subsection{Что \texttt{auto} делает с references и const}

Обычный \texttt{auto} deduction во многом похож на template type deduction.

\begin{cppcode}
const int value = 10;
const int& ref = value;

auto a = ref;        // int
const auto b = ref;  // const int
const auto& c = ref; // const int&
\end{cppcode}

Top-level const/reference у initializer не обязательно сохраняются простым
\texttt{auto}. Если semantics требует reference, нужно написать
\texttt{auto\&} или \texttt{const auto\&} явно.

\begin{warningbox}[Auto не должен скрывать ownership или lifetime]
Если функция возвращает reference, запись \texttt{auto x = f();} может создать
copy вместо alias. Если нужен borrow, пишите \texttt{auto\&} или
\texttt{const auto\&} и проверяйте lifetime.
\end{warningbox}

\subsection{\texttt{auto\&}, \texttt{const auto\&}, \texttt{auto\&\&}}

Практическая таблица:

\begin{center}
\begin{tabularx}{\textwidth}{L{0.22\textwidth}YY}
\toprule
Форма & Идея & Обычный смысл \\
\midrule
\texttt{auto x} & новое value & copy/move result \\
\texttt{auto\& x} & mutable borrow & изменить существующий object \\
\texttt{const auto\& x} & read-only borrow & избежать copy \\
\texttt{auto\&\& x} & deduction с value category & generic/forwarding context \\
\bottomrule
\end{tabularx}
\end{center}

После блока~\ref{sec:lecture-12} последняя строка уже не выглядит магией:
\texttt{auto\&\&} следует deduction и reference-collapsing rules.

\subsection{\texttt{decltype}: type expression без вычисления}

\texttt{decltype(expression)} позволяет получить type по правилам языка, не
выполняя expression ради side effects.

\begin{cppcode}
int value = 7;
decltype(value) copy = 9;   // int
decltype((value)) alias = value; // int&
\end{cppcode}

Разница двойных скобок принципиальна. Для unparenthesized id-expression
\texttt{decltype(name)} обычно даёт declared type entity. Для общего expression
результат зависит от value category:
lvalue expression $\rightarrow$ \texttt{T\&}; xvalue expression $\rightarrow$ \texttt{T\&\&}; prvalue expression $\rightarrow$ \texttt{T}.

\begin{mistakebox}[Скобки вокруг имени меняют результат]
\texttt{decltype(x)} и \texttt{decltype((x))} могут быть разными types.
Вторая запись рассматривает \texttt{(x)} как обычное expression и учитывает его
value category.
\end{mistakebox}

\subsection{\texttt{decltype(auto)}}

Return type \texttt{decltype(auto)} применяет правила \texttt{decltype} к
return expression:

\begin{cppcode}
decltype(auto) at(std::vector<int>& values, std::size_t i) {
    return (values[i]);
}
\end{cppcode}

Скобки здесь намеренно сохраняют reference. Это мощный, но опасный инструмент:
если вернуть reference на object с закончившимся lifetime, deduction не спасёт.

Для простого API явный return type часто читабельнее.

\subsection{Trailing return type}

Иногда return type зависит от parameters:

\begin{cppcode}
template <typename A, typename B>
auto add(const A& a, const B& b) -> decltype(a + b) {
    return a + b;
}
\end{cppcode}

В C++14+ многие функции могут просто использовать deduced \texttt{auto} return
type, но trailing form остаётся полезной, когда return type нужно явно выразить
через parameter expressions.

\subsection{Structured bindings}

C++17 позволяет разложить tuple-like object или подходящую structure:

\begin{cppcode}
std::map<std::string, int> scores{
    {"templates", 9},
    {"move", 10}
};

for (const auto& [topic, score] : scores) {
    std::cout << topic << ": " << score << '\n';
}
\end{cppcode}

Имена \texttt{topic} и \texttt{score} связываются с компонентами элемента.
Важно выбрать форму binding:
\texttt{auto [a,b]} и \texttt{auto\& [a,b]} имеют разные copy/reference
semantics.

\begin{warningbox}[Не забывайте \& в range-for]
Если нужно изменять members элемента container, простое
\texttt{auto [key, value]} может работать с отдельным binding object/copy,
а \texttt{auto\& [key, value]} выражает работу с самим элементом там, где это
разрешено типом.
\end{warningbox}

\subsection{CTAD: class template argument deduction}

До C++17 template arguments class часто писали явно:

\begin{cppcode}
std::pair<std::string, int> item{"templates", 9};
\end{cppcode}

В C++17:

\begin{cppcode}
std::pair item{std::string{"templates"}, 9};
\end{cppcode}

Compiler использует deduction guides и constructors, чтобы вывести template
arguments. Это называется \term{CTAD}.

\begin{standardbox}[CTAD не делает class template «нетипизированным»]
После deduction объект всё равно имеет конкретный specialization, например
\texttt{std::pair<std::string, int>}. Deduction происходит на этапе
компиляции.
\end{standardbox}

Для user-defined class templates compiler может использовать implicit deduction
guides, а programmer при необходимости может объявлять свои. В этом блоке
достаточно понимать сам механизм; сложные guides оставим библиотечному design.

\subsection{Callable object: общий термин}

В C++ «что-то можно вызвать» не означает только обычную функцию. Callable может
быть:
function; function pointer; class object с \texttt{operator()}; lambda closure object; wrapper вроде \texttt{std::function}.

Templates особенно хорошо работают с callables, потому что конкретный type
известен compiler.

\subsection{Function object}

До lambda поведение часто упаковывали в class:

\begin{cppcode}
struct Threshold {
    int limit;

    bool operator()(int value) const {
        return value >= limit;
    }
};
\end{cppcode}

Это обычный object со state и member function \texttt{operator()}. Lambda
делает такой локальный pattern намного компактнее.

\subsection{Минимальная lambda}

\begin{cppcode}
auto is_even = [](int value) {
    return value % 2 == 0;
};

bool result = is_even(42);
\end{cppcode}

Синтаксис:
1) \texttt{[]} --- capture list; 2) \texttt{(int value)} --- parameters; 3) optional specifiers; 4) body в фигурных скобках.

\subsection{Что создаёт compiler}

Каждая lambda expression создаёт prvalue уникального unnamed
\term{closure type}. Концептуально можно представить:

\begin{cppcode}
class SomeUniqueClosure {
public:
    bool operator()(int value) const {
        return value % 2 == 0;
    }
};
\end{cppcode}

Это не буквальный required layout и не имя, которым можно воспользоваться.
Модель нужна, чтобы понять:
lambda variable --- object, его можно копировать/перемещать согласно свойствам
closure type, а вызов похож на вызов \texttt{operator()}.

\subsection{Capture by value}

\begin{cppcode}
int threshold = 10;

auto enough = [threshold](int value) {
    return value >= threshold;
};
\end{cppcode}

Closure хранит захваченное значение. Дальнейшее изменение внешнего
\texttt{threshold} не обязано менять сохранённую copy.

Value capture особенно удобен, когда callback должен пережить локальную
variable, а захваченный object сам безопасно копируется.

\subsection{Capture by reference}

\begin{cppcode}
int accepted = 0;

auto count = [&accepted](int value) {
    if (value > 0) {
        ++accepted;
    }
};
\end{cppcode}

Closure хранит возможность обратиться к внешнему object, но не продлевает его
lifetime.

\begin{warningbox}[Reference capture --- borrowing contract]
Если lambda сохраняется дольше захваченной local variable, последующий вызов
может обратиться к object после окончания lifetime. Это тот же класс проблемы,
что dangling reference/pointer.
\end{warningbox}

Поэтому callback registry, который хранит lambda, требует особенно внимательно
читать capture list.

\subsection{Default capture}

\texttt{[=]} задаёт default capture by value, \texttt{[\&]} --- by reference.

Коротко, но большие default captures могут скрывать реальные зависимости.
В callback, который живёт долго, explicit capture list часто лучше:

\begin{cppcode}
auto callback = [prefix, &counter](const std::string& value) {
    ++counter;
    std::cout << prefix << value << '\n';
};
\end{cppcode}

По коду видно, что именно closure хранит/заимствует.

\subsection{\texttt{mutable} lambda}

По умолчанию call operator non-mutable lambda, которая захватывает by value,
const относительно сохранённого state.

\begin{cppcode}
auto next = [counter = 0]() mutable {
    return ++counter;
};
\end{cppcode}

\texttt{mutable} разрешает изменять copy внутри closure. Внешний object, если он
вообще существовал отдельно, этим не изменяется.

\subsection{Init-capture}

C++14 добавил возможность создать capture прямо в списке:

\begin{cppcode}
auto callback =
    [name = std::string{"audit"}]() {
        return name.size();
    };
\end{cppcode}

Это удобно для rename, conversion и move-capture:

\begin{cppcode}
auto owner = std::make_unique<Job>();

auto task =
    [job = std::move(owner)]() {
        job->run();
    };
\end{cppcode}

После такого capture closure может стать move-only, потому что содержит
\texttt{unique\_ptr}. Это важно для \texttt{std::function} в C++17:
он требует CopyConstructible target.

\subsection{Generic lambda}

C++14 разрешает \texttt{auto} в lambda parameters:

\begin{cppcode}
auto greater = [](const auto& a, const auto& b) {
    return a > b;
};
\end{cppcode}

Концептуально call operator closure --- function template. Один closure
object может работать с несколькими types, если body корректен для каждого
instantiation.

Это естественное продолжение templates из блока~\ref{sec:lecture-09}.

\subsection{Lambda return type}

Часто return type выводится:

\begin{cppcode}
auto square = [](int x) {
    return x * x;
};
\end{cppcode}

Если нужна явность:

\begin{cppcode}
auto parse = [](const std::string& text) -> std::optional<int> {
    // ...
};
\end{cppcode}

Return expressions должны удовлетворять правилам deduction. Если branches
возвращают несовместимые types, compiler сообщит ошибку.

\subsection{Lambdas и algorithms}

Следующий блок подробно разберёт STL algorithms. Уже сейчас видно главное
применение:

\begin{cppcode}
auto it = std::find_if(
    blocks.begin(),
    blocks.end(),
    [](const CourseBlock& block) {
        return block.ready;
    }
);
\end{cppcode}

Algorithm отвечает за traversal, lambda --- только за policy/condition.
Это пример разделения механизма и поведения без class hierarchy.

\subsection{Когда template callable лучше}

Generic function:

\begin{cppcode}
template <typename Predicate>
bool any_match(const std::vector<CourseBlock>& blocks,
               Predicate predicate);
\end{cppcode}

получает concrete callable type в каждом instantiation. Compiler может inline
call и оптимизировать через границу callback.

Плюсы:
нет обязательного type-erasure wrapper; callable может быть move-only; concrete type доступен compiler.

Минусы:
implementation обычно должна быть видна в header; разные callable types создают разные instantiations; heterogeneous callbacks нельзя просто хранить в одном \texttt{std::vector<Predicate>}.

\subsection{\texttt{std::function}: единый runtime callable type}

\texttt{std::function<R(Args...)>} из \headerfile{functional} может хранить
разные callable objects с совместимой сигнатурой:

\begin{cppcode}
std::function<void(const Event&)> callback;

callback = free_function;
callback = SomeFunctor{};
callback = [prefix](const Event& event) {
    std::cout << prefix << event.name;
};
\end{cppcode}

Внешний type один и тот же, хотя внутренний callable type различается. Это
\term{type erasure}.

\subsection{Зачем нужна type erasure}

Представим registry:

\begin{cppcode}
using Callback = std::function<void(const Event&)>;

std::vector<Callback> callbacks;
\end{cppcode}

Теперь callbacks можно выбирать во время выполнения и хранить в одном
container. Template parameter одного compile-time call site этого сам по себе
не даёт.

Именно runtime heterogeneity --- основное основание для
\texttt{std::function}, а не желание «завернуть любую lambda».

\subsection{Стоимость \texttt{std::function}}

Стандарт задаёт semantics, но не требует конкретный layout или конкретную
small-buffer optimization.

Практически wrapper может вносить:
indirect call; дополнительное состояние для type erasure; иногда dynamic allocation для target; невозможность некоторых optimizations по сравнению с полностью известным concrete callable type.

\begin{performancebox}[Не говорить «std::function всегда выделяет память»]
Реализации часто используют small-object optimization, но размер буфера и факт
allocation для конкретного closure не являются универсальным обещанием,
на которое должен опираться переносимый код.
\end{performancebox}

Если callback только немедленно вызывается, template parameter часто проще и
дешевле. Если callback нужно хранить heterogeneously, \texttt{std::function}
может быть именно тем abstraction, который нужен.

\subsection{Пустой \texttt{std::function}}

Wrapper может не содержать callable:

\begin{cppcode}
std::function<void()> callback;

if (callback) {
    callback();
}
\end{cppcode}

Вызов пустого \texttt{std::function} бросает
\texttt{std::bad\_function\_call}. Обычно registry проверяет наличие или
гарантирует invariant, что slot заполнен.

\subsection{Copyability target в C++17}

Для конструкции \texttt{std::function} из callable target в C++17 target должен
удовлетворять требованиям copyability.

Поэтому:

\begin{cppcode}
auto owner = std::make_unique<Job>();

auto move_only =
    [job = std::move(owner)]() {
        job->run();
    };

// std::function<void()> f = std::move(move_only);
// для move-only closure это не подходит в C++17
\end{cppcode}

Эксперимент проверяет это настоящей compile failure.

\subsection{\texttt{std::optional<T>}: отсутствие как часть типа}

До C++17 функции часто кодировали «не найдено» через:
magic value \texttt{-1}; пустую строку; nullable pointer; отдельный bool + output parameter.

Если функция концептуально возвращает value \texttt{T}, но value может
отсутствовать, \texttt{std::optional<T>} делает это явным:

\begin{cppcode}
std::optional<CourseBlock> find_block(int number);
\end{cppcode}

Result либо содержит \texttt{CourseBlock}, либо находится в disengaged state.

\subsection{Создание optional}

\begin{cppcode}
std::optional<int> a = 42;
std::optional<int> b = std::nullopt;
auto c = std::make_optional<std::string>("ready");
\end{cppcode}

Для функции удобно:

\begin{cppcode}
std::optional<int> find_score(const std::string& topic) {
    if (topic == "templates") {
        return 9;
    }

    return std::nullopt;
}
\end{cppcode}

\subsection{Проверка optional}

\begin{cppcode}
if (auto score = find_score("templates")) {
    std::cout << *score;
}
\end{cppcode}

Можно использовать:
contextual bool conversion; \texttt{has\_value()}; \texttt{operator*} / \texttt{operator->} после доказанного наличия; \texttt{value()}; \texttt{value\_or(default)}.

\texttt{value()} у пустого optional бросает
\texttt{std::bad\_optional\_access}. Разыменовывать пустой optional через
\texttt{operator*} нельзя.

\subsection{Optional не означает error details}

\texttt{optional<T>} отвечает только на вопрос «есть T или нет». Он не сообщает,
\emph{почему} значения нет.

Если caller должен различать invalid input, I/O failure и not-found, нужны
другие error contracts: status object, exception или более богатый result type.

Это связывает тему с блоком~\ref{sec:lecture-10}: abstraction выбирается по
информации, которую обязан получить caller.

\subsection{Optional и pointers}

Если функция возвращает non-owning view на существующий object, nullable pointer
может быть естественнее:

\begin{cppcode}
const CourseBlock* find_block_view(int number);
\end{cppcode}

Если функция возвращает самостоятельное value, которое может отсутствовать:

\begin{cppcode}
std::optional<CourseBlock> find_block_copy(int number);
\end{cppcode}

Выбор зависит не от моды, а от ownership/lifetime contract.

\subsection{Optional reference отсутствует}

В C++17 стандартный \texttt{std::optional<T\&>} не поддерживается.
Для optional borrow обычно используют pointer или
\texttt{std::reference\_wrapper<T>} внутри optional, если это действительно
улучшает design.

Не пытайтесь превращать optional в универсальный контейнер «для всего».

\subsection{Прикладной query API}

Соберём темы вместе:

\begin{cppcode}
template <typename Predicate>
std::optional<CourseBlock>
find_first(const std::vector<CourseBlock>& blocks,
           Predicate predicate) {
    for (const auto& block : blocks) {
        if (predicate(block)) {
            return block;
        }
    }

    return std::nullopt;
}
\end{cppcode}

Caller:

\begin{cppcode}
const int min_examples = 2;

auto result = find_first(
    blocks,
    [min_examples](const CourseBlock& block) {
        return block.examples >= min_examples;
    }
);

if (result) {
    std::cout << result->slug;
}
\end{cppcode}

Здесь:
template выводит concrete predicate type; lambda хранит threshold by value; \texttt{const auto\&} внутри query не копирует элементы; \texttt{optional<CourseBlock>} явно выражает отсутствие; caller не использует magic index \texttt{-1}.

\subsection{Structured binding в прикладном API}

Если функция возвращает несколько связанных значений:

\begin{cppcode}
auto summary = std::pair{ready_count, total_count};
const auto [ready, total] = summary;
\end{cppcode}

CTAD выводит template arguments \texttt{pair}, structured binding даёт имена
двум components. Но если у пары есть domain meaning, отдельный named struct
может быть читабельнее.

\begin{whybox}[Современный синтаксис не отменяет моделирование]
\texttt{auto}, CTAD и structured bindings уменьшают шум. Они не должны
превращать программу в набор безымянных tuple/pair, если предметная область
заслуживает собственного типа.
\end{whybox}

\subsection{Callback registry как runtime boundary}

Registry хранит разные callbacks:

\begin{cppcode}
class Registry {
public:
    using Callback = std::function<void(const Event&)>;

    void subscribe(Callback callback) {
        callbacks_.push_back(std::move(callback));
    }

    void publish(const Event& event) const {
        for (const auto& callback : callbacks_) {
            callback(event);
        }
    }

private:
    std::vector<Callback> callbacks_;
};
\end{cppcode}

Здесь \texttt{std::function} оправдан: callbacks имеют разные closure types,
создаются в разных местах и должны храниться вместе после завершения
\texttt{subscribe}.

\subsection{Capture lifetime в registry}

Опасный вариант:

\begin{cppcode}
Registry registry;

{
    int local = 0;
    registry.subscribe([&local](const Event&) {
        ++local;
    });
} // local уничтожен

// registry.publish(...); // dangling reference внутри closure
\end{cppcode}

\texttt{std::function} корректно владеет closure object, но closure object может
сам содержать dangling reference. Type erasure не исправляет capture lifetime.

Безопасные варианты:
capture independent state by value; захватывать \texttt{shared\_ptr}, если callback действительно должен продлевать shared lifetime; использовать \texttt{weak\_ptr}, если callback только наблюдает owner; гарантировать внешний lifetime reference capture явным design contract.

\subsection{Mutable state внутри callback}

Init-capture позволяет callback иметь собственное состояние:

\begin{cppcode}
registry.subscribe(
    [calls = 0](const Event&) mutable {
        ++calls;
        std::cout << calls << '\n';
    }
);
\end{cppcode}

Каждая copy closure получает собственную copy state. Если wrapper/registry
копирует callback, это может иметь значение. Нельзя автоматически считать,
что mutable lambda --- «один глобальный счётчик».

\subsection{Lambda и \texttt{this}}

В member function lambda может обращаться к object через capture \texttt{this}.
Но если callback переживёт object, raw \texttt{this} станет dangling.

Для long-lived callback ownership model нужно проектировать явно. Например,
shared-owned object иногда передаёт \texttt{weak\_ptr} в callback и перед
использованием вызывает \texttt{lock()}.

Так блоки 11 и 13 соединяются напрямую.

\subsection{Когда function pointer всё ещё достаточен}

Captureless lambda может преобразовываться к совместимому function pointer.
Обычный function pointer остаётся хорошим инструментом для C API или API, где
state не нужен.

\texttt{std::function} не нужно использовать автоматически в каждом месте,
где параметр можно вызвать.

\subsection{Generic lambda или std::function?}

Сравнение:

\begin{center}
\begin{tabularx}{\textwidth}{L{0.25\textwidth}YY}
\toprule
Свойство & Template/generic callable & \texttt{std::function} \\
\midrule
Concrete type известен compiler & да & скрыт type erasure \\
Heterogeneous storage & не напрямую & да \\
Move-only callable в C++17 & да, если API поддерживает move & target для \texttt{std::function} должен быть copyable \\
Inlining opportunities & обычно лучше видны compiler & indirect boundary может мешать \\
Stable runtime wrapper type & нет одного общего type & да \\
\bottomrule
\end{tabularx}
\end{center}

\subsection{Эксперимент 1: course\_query\_pipeline}

Каталог:
\filepath{examples/13\_type\_deduction\_lambdas\_and\_function/01\_course\_query\_pipeline}.

Эксперимент строит маленький query API для учебных blocks:
1) \texttt{auto}, \texttt{decltype} и \texttt{static\_assert} проверяют реальные deduced types; 2) structured bindings разбирают summary; 3) CTAD создаёт \texttt{std::pair} без явных template arguments; 4) generic lambda работает с несколькими numeric types; 5) predicate lambda фильтрует blocks; 6) \texttt{std::optional<CourseBlock>} представляет found/not-found без magic index.

Отдельный negative compile file действительно использует local variable в
lambda без capture и должен быть отвергнут compiler.

\begin{shellcode}
cd examples/13_type_deduction_lambdas_and_function/01_course_query_pipeline
./run_checks.sh
\end{shellcode}

\subsection{Эксперимент 2: callback\_registry\_lab}

Каталог:
\filepath{examples/13\_type\_deduction\_lambdas\_and\_function/02\_callback\_registry\_lab}.

Registry хранит \texttt{std::function<void(const Event\&)>} и принимает:
free function; function object; lambda с value capture; mutable lambda с собственным state; lambda, наблюдающую shared-owned object через \texttt{weak\_ptr}.

Проверяется также empty \texttt{std::function} и
\texttt{std::bad\_function\_call}. Negative compile case показывает ограничение
C++17: move-only closure с \texttt{unique\_ptr} нельзя положить в
\texttt{std::function}.

\begin{shellcode}
cd examples/13_type_deduction_lambdas_and_function/02_callback_registry_lab
./run_checks.sh
\end{shellcode}

\subsection{Стоимость и границы abstraction}

\texttt{auto}, \texttt{decltype}, structured bindings и CTAD в основном влияют
на compile-time type deduction и сами по себе не добавляют отдельный runtime
object «для вывода типа».

Lambda closure --- настоящий object. Его size зависит от captures и
implementation details/alignment. Captureless closure не обязана иметь size
zero.

Template callable часто позволяет compiler видеть concrete call path.
\texttt{std::function} покупает runtime uniformity ценой type-erasure layer и
возможного allocation/indirect call.

\texttt{optional<T>} обычно должен хранить состояние engaged/disengaged плюс
место для \texttt{T}, но конкретный layout не гарантируется как простая формула
\texttt{sizeof(T)+1}.

\begin{performancebox}[Не угадываем стоимость по синтаксису]
Короткая lambda не обязательно «быстрее функции», а \texttt{std::function} не
обязательно делает allocation. Если callback path важен для performance,
измеряют реальную программу и конкретные target/compiler settings.
\end{performancebox}

\subsection{Итоговый checklist}

Когда видите длинный type:
если initializer однозначно задаёт type --- рассмотрите \texttt{auto}; если важны reference/const semantics --- напишите их явно; если нужен type expression --- \texttt{decltype}; для tuple-like decomposition --- structured bindings; CTAD используйте, когда deduction делает type очевиднее, а не загадочнее.

Когда нужно передать поведение:
local behavior --- lambda; generic immediate call --- template callable; runtime heterogeneous storage --- \texttt{std::function}; capture lifetime проверяется отдельно от wrapper.

Когда результат может отсутствовать:
standalone value --- \texttt{std::optional<T>}; non-owning optional view --- часто pointer; если нужны причины ошибки --- более богатый error contract.

\begin{importantbox}[Главный вывод]
Современный C++ позволяет писать меньше типов вручную, но правильный design всё
равно строится вокруг точных contracts. \texttt{auto} не отменяет type,
lambda не отменяет lifetime, \texttt{std::function} не отменяет стоимость
runtime abstraction, а \texttt{optional} не отменяет проектирование ошибок.
\end{importantbox}

Следующий блок переходит к STL как к единой системе containers, iterators и
algorithms. Там callable objects этой лекции станут обычными predicates,
comparators и operations для standard algorithms.
